\section{讨论段落判别样例}

\subsection{样例 A：无证据的机制堆叠}
引入辅助损失后，宏平均 F1 提高了 3.2 个百分点。这一增益可能源于辅助损失平滑了优化
景观，也可能源于编码器学到了更稳定的特征，还可能是解码器的校准能力有所改善。
现有比较没有测量上述任何机制，也无法验证其中哪一项导致了该增益。

\subsection{样例 B：单一解释有局部证据且使用合理限定}
误差下降可能来自预测方差减小：在 10 个随机种子上，方差降低了 18\%；移除集成后，
方差恢复到原有水平。

\subsection{样例 C：多个假设均有证据和区分性检验}
两项机制共同带来了性能提升。类别均衡增强使少数类召回率提高 6.1 个百分点，且没有改变
校准误差；温度缩放则把期望校准误差从 0.082 降至 0.041，且没有改变召回率。独立消融
区分了两项机制各自的贡献。

\subsection{样例 D：明确承认机制尚未确定}
引入辅助损失后，宏平均 F1 提高了 3.2 个百分点。当前实验设计无法确定这一差异的形成
机制。

\subsection{样例 E：受控实验支持强机制措辞}
在随机化消融中，关闭检索模块后 3.2 个百分点的增益完全消失，其余组件与随机种子均保持
不变。在该基准范围内，检索模块导致了观察到的性能提升。
